% 由 scripts/publication/tables.py 自动生成，请勿手改。
% 需要宏包：booktabs（\usepackage{booktabs}）
\begin{table}[htbp]
  \centering
  \caption{512 通道交叉损耗模型的敏感性（固定几何复算）}
  \label{tab:crossing-sensitivity}
  \begin{tabular}{llrr}
    \toprule
    口径 & 说明 & 平均损耗 (dB) & 最大损耗 (dB) \\
    \midrule
    A & 全部按 90°（0.05 dB/30），忽略角度依赖 & 5.4758 & 6.5224 \\
    B & 图 3-12 数字化 + 90° 锚定（项目主口径） & 5.5147 & 6.5761 \\
    C & 图 3-12 原始值（不锚定） & 5.5745 & 6.6659 \\
    \bottomrule
  \end{tabular}
  \par\vspace{2pt}
  \begin{flushleft}\footnotesize
  三种口径均为固定几何复算（不改布线）：非交叉部分相同，只替换交叉项；逐路交叉条目按每根波导各计一次。 \\
  与 OpticalWaveguideRouter2D/docs/loss\_model\_reconstruction\_report.md 第 7 节的记录值（均值 5.4758/5.5147/5.5745 与最大 6.5224/6.5761/6.6659）一致（±0.002 dB，差异来自文档记录的四位舍入）。 \\
  主口径（B）即全部复现图的损耗口径。 \\
  \end{flushleft}
\end{table}
